\section{Privacy Analysis and Budget Allocation}
\label{sec:analysis}
We first prove the privacy guarantee for the PrivTDD protocol and then derive an allocation of its reporting budgets.
% First establish privacy for the complete protocol, then derive report-budget shares from the expected error bound for the released histogram.

\subsection{Privacy}
\begin{theorem}[Privacy of the complete protocol]\label{thm:pg-private}
The complete PrivTDD transcript satisfies pure $\varepsilon$-edge privacy when
\[
  2\varepsilon_d+\varepsilon_r+2\varepsilon_\ell\leq\varepsilon.
\]
Here $\varepsilon_d=\varepsilon_{\mathrm{deg}}$, $\varepsilon_r=\varepsilon_{\mathrm{rr}}$, and $\varepsilon_\ell=\varepsilon_{\mathrm{rel}}$ are the degree, RR, and relational-report budgets in Eq.~\eqref{eq:privacy_allocation}.
\end{theorem}

\begin{proof}
Changing an undirected edge incurs $2\varepsilon_d$ for its endpoint degree reports and $\varepsilon_r$ for its assigned RR report.
Fix the graph-independent public neighbor order and the first-round transcript, including the caps and shared RR table.
For an endpoint with cap $c>1$, growth of its retained set adds at most $c-1$ centered terms of magnitude at most one.
Replacing one retained neighbor preserves the pair count, so the centered-count difference is a sum of at most $c-1$ differences between RR bits, each of magnitude at most one.
Deletion is symmetric, and the centered count $B_v-p_0\widetilde{W}_v$ is zero for $c\leq1$.
Thus the conditional sensitivity is at most $(c-1)_+$.
At Laplace scale $(c-1)_+/\varepsilon_\ell$, the second round costs at most $\varepsilon_\ell$ per endpoint.
Other nodes' conditional report laws are unchanged.
Composition gives the stated budget, and estimation and the safeguard are post-processing.
\end{proof}

\subsection{Budget Allocation}
\label{sec:analysis:allocation}
\label{sec:analysis:utility}

A practical challenge is deciding how to allocate a fixed privacy budget among PrivTDD's three reporting mechanisms before data collection.
The degree budget controls noise in $C_v$, which affects both neighbor truncation and the estimated wedge-count distribution.
The relational-report budget controls the Laplace noise added to the centered count.
The RR budget affects how well randomized adjacency bits distinguish edges from nonedges and how much that same Laplace noise is amplified during debiasing.
Theorem~\ref{thm:pg-private} limits the weighted sum of these budgets. Giving more to one report leaves less for the others, and the server lacks the true TDD needed to compare the resulting errors.
% Budgets are chosen before collection. The degree budget affects truncation and wedge counts; the relational budget controls Laplace noise.
% The RR budget affects edge distinguishability and noise amplification during debiasing. The true histogram is unavailable for choosing shares.


Our idea is to derive an upper bound on the expected EMD of the released histogram and minimize it before collection.
Let $h=H/n$, $h_A=H_A/n$, and $\widehat h=\widehat H/n$ denote the normalized true, reference, and released histograms. Thus $\operatorname{EMD}(\widehat H,H)=W_1(\widehat h,h)$.
Phase~4 forms $H_A$ from binned relational reports and estimated wedge counts. For any graph on $n$ nodes, the reporting mechanisms give an error bound for each input in terms of the candidate budgets.
The release rule in Eq.~\eqref{eq:pg-guard} keeps $W_1(\widehat h,h_A)\leq\tau D$, with $D=D_{\max}=n-1$ and fixed public $\tau>0$.
Consequently, an expected error bound $B_A$ for $h_A$ gives $B_A+\tau D$ for $\widehat h$. $B_A$ varies across allocations.
% Direct binning and wedge-count estimation contribute two reference-histogram errors; Phase 4 relates both to the released histogram.
% Histograms have total count n and are normalized only for distances. The fixed radius means allocation compares only B_A.

\begin{lemma}[Error of the reference histogram]\label{lem:pg-reference}
Suppose Phase~1 uses all retained neighbor pairs and fits the degree proportions $\widehat\pi$ by exact maximum likelihood on $\{0,\ldots,D\}$.
Let $p=e^{\varepsilon_r}/(1+e^{\varepsilon_r})$ and $p_0=1-p$ be the RR probabilities of reporting one for an edge and a nonedge, as in Section~\ref{sec:private_collection}. Set $\gamma=p-p_0$ and $W_{\max}=\binom D2$.
Averaging over protocol randomness for any fixed $G$ gives $\mathbb E_G W_1(h_A,h)\leq B_A$, where
\begin{equation}
\begin{aligned}
B_A={}&\frac{\sqrt{pp_0 W_{\max}}}{\gamma}
       +\frac{D-1}{\gamma\varepsilon_\ell}
       +\frac{D-1}{2(1-e^{-\varepsilon_d})}+\frac12\\
     &+(D-1)\bigl[\coth(\varepsilon_d/2)+1/\varepsilon_d+1\bigr].
\end{aligned}
\label{eq:pg-anchor-bound}
\end{equation}
\end{lemma}
% This bound assumes full reporting and exact full-degree-domain maximum likelihood. The first line accounts for RR noise, Laplace noise, truncation, and rounding; the second accounts for wedge-count error.
\begin{proof}
Pair each binned relational report with its true triangle degree $t_v$. Given $C_v$ and the fixed graph, the RR count has standard deviation at most $\sqrt{pp_0W_{\max}}$ and the expected absolute Laplace noise is at most $(D-1)/\varepsilon_\ell$. Debiasing divides both by $\gamma$.
Each omitted neighbor costs at most $D-1$ triangles. With $a=\varepsilon_d$, Eq.~\eqref{eq:private_degree} gives $\mathbb E_G(d_v-C_v)_+=\tfrac12\sum_{k=0}^{d_v-1}e^{-ak}\leq[2(1-e^{-a})]^{-1}$.
Clipping cannot increase error and rounding adds at most $1/2$. This nodewise coupling allows shared RR bits.
% Pair each report with its true triangle degree to accumulate error; independence across nodes' RR reports is not required.

For the wedge counts, write $K=\mathcal K_{\mathrm{deg}}$ for the degree-report channel in Eq.~\eqref{eq:degree_distribution}. Set $f(c)=\sum_d\widehat\pi(d)K(c\mid d)$ and $g_j=n^{-1}\sum_v K(C_v\mid j)/f(C_v)$.
At the maximum, the log-likelihood derivative toward unit mass at $j$ is $n(g_j-1)\leq0$. Since $\sum_j\widehat\pi(j)g_j=1$, the mean degree posterior is $\bar q(d):=n^{-1}\sum_v q_v(d)=\widehat\pi(d)g_d=\widehat\pi(d)$.
For the floored and clipped Laplace channel,
\[
 |d-c|K(c\mid d)\leq\sum_{j=0}^D e^{-a|j-d|}K(c\mid j).
\]
For each integer $j$ from $d$ toward $c$, excluding $c$, the Laplace channel gives $K(c\mid j)=e^{a|j-d|}K(c\mid d)$.
Set $c=C_v$, multiply by $\widehat\pi(d)/(n f(C_v))$, and sum over $v,d$.
Using $g_j\leq1$ gives $n^{-1}\sum_{v,d}q_v(d)|d-C_v|\leq\sum_{k\in\mathbb Z} e^{-a|k|}=\coth(a/2)$.
Coupling each cap to its true degree adds expected cost at most $1/a+1$. The map $w(d)=\binom d2$ is $(D-1)$-Lipschitz, giving the last term of $B_A$.
% Full-domain first-order conditions equate the mean degree posterior and fitted degree law. Two couplings bound posterior-to-cap and cap-to-true-degree distances.
% The wedge-count map amplifies this distance by at most D-1.

Let $F_G$ and $F_{\mathrm{cap},G}$ be the true triangle-degree and wedge-count CDFs. Since $t_v\leq w(d_v)$, we have $F_G\geq F_{\mathrm{cap},G}$ pointwise.
The maximum in Eq.~\eqref{eq:pg-reference-cdf} gives
$|F_A-F_G|\leq|F_{\mathrm{dir}}-F_G|+|F_{\mathrm{cap}}-F_{\mathrm{cap},G}|$.
Integrating proves the claim~\cite{vallender1974distance}.
\end{proof}

\begin{theorem}[Minimax EMD for TDD]\label{thm:pg-risk}
Under the assumptions of Lemma~\ref{lem:pg-reference}, the first inequality holds for every feasible positive budget allocation. The second line gives the minimax order for fixed $\varepsilon>0$.
\begin{equation}
\begin{aligned}
 \sup_G\mathbb E_G\operatorname{EMD}(\widehat H,H)
     &\leq B_A+\tau D,\\
 \inf_\Pi\sup_G\mathbb E_G\operatorname{EMD}(\Pi(G),H)
     &=\Theta_{\varepsilon}(n).
\end{aligned}
\label{eq:pg-minimax}
\end{equation}
Here $G$ ranges over simple graphs on the public vertex set and $\Pi$ over local protocols with the same edge privacy budget that release a histogram $\Pi(G)$ of total count $n$.
For fixed $\varepsilon,\tau>0$ and budget shares bounded away from 0, $B_A+\tau D=O(n)$. Unnormalized transport cost is $n\operatorname{EMD}$ with minimax order $\Theta_\varepsilon(n^2)$.
\end{theorem}
% The finite-sample upper bound holds for every feasible positive split; the lower bound gives worst-case growth in n at fixed total budget.
\begin{proof}
By the mixture in Eq.~\eqref{eq:pg-guard}, the one-dimensional CDF formula gives $W_1(\widehat h,h_A)=\alpha\Delta\leq\tau D$. The triangle inequality and Lemma~\ref{lem:pg-reference} give the upper bound. For the lower bound, set $m=\lfloor(n-2)/4\rfloor$ for $n\geq6$ and form a clique on $2m+2$ vertices. Join each new vertex pair $i$, $1\leq i\leq m$, to the first $2i$ clique vertices. An independent fair bit $X_i$ determines the edge within pair $i$. Leave the remaining vertices isolated.
Each vertex of pair $i$ has triangle degree $a_i=\binom{2i}{2}$ or $a_i+2i$. The intervals $[a_i,a_i+2i)$ are disjoint and all clique vertices have at least $a_m+2m$ triangles.
At each of the $2i$ integer cutoffs in interval $i$, the cumulative count therefore takes two values differing by two according to $X_i$.
Edge privacy bounds the posterior odds of $X_i$ given the transcript between $e^{-\varepsilon}$ and $e^\varepsilon$, giving expected absolute count error at least $2/(1+e^\varepsilon)$ per cutoff.
Summing over cutoffs and dividing by $n$ gives average risk at least $2m(m+1)/[n(1+e^\varepsilon)]=\Omega_\varepsilon(n)$. Worst-case risk is at least as large.
\end{proof}

\emph{Approximate degree fits.}
For an approximate fit, let $r_{\mathrm{opt}}=\max_{0\leq j\leq D}g_j-1$ with the scores defined above.
The total variation distance between $\bar q$ and $\widehat\pi$ is at most $r_{\mathrm{opt}}$ since $\bar q(d)=\widehat\pi(d)g_d$. The posterior-to-cap coupling now costs at most $(1+r_{\mathrm{opt}})\coth(\varepsilon_d/2)$.
Thus the upper bound increases by at most $(D-1)[D+\coth(\varepsilon_d/2)]\mathbb E_G r_{\mathrm{opt}}$, using degree-domain diameter $D$.
For fixed $\varepsilon,\tau>0$ and budget shares bounded away from zero, the $O(n)$ rate holds if $\sup_G\mathbb E_G r_{\mathrm{opt}}=O(1/D)$. Even for a restricted grid, this residual uses scores at every degree.
% Restricted-grid or early-stopped fits require a full-domain residual check. The allocation below optimizes the public exact-fit bound before collection.

\noindent\textbf{Allocation scheme.}
We choose the budgets using the explicit full-domain, exact-fit bound $B_A+\tau D$ in Eq.~\eqref{eq:pg-minimax}. Its dependence on each budget distinguishes splits with the same $O(n)$ rate.
With the radius fixed, $B_A$ decreases in each of the three budgets. The minimizing split therefore exhausts the budget.
Writing $x=\varepsilon_d$ and $y=\varepsilon_r$ gives $\varepsilon_\ell=(\varepsilon-2x-y)/2$.
% Fix public inputs before collection and exhaust the privacy budget. Eliminating the relational-report budget leaves two optimization variables.

\begin{proposition}[Bound-optimal allocation]\label{prop:pg-allocation}
The allocation minimizing $B_A+\tau D$ is determined by the unique solution $(x^\star,y^\star)$ of
\begin{equation}
\begin{aligned}
\min_{\substack{x,y>0\\2x+y<\varepsilon}}\quad
 &\frac{\sqrt{D/[2(D-1)]}}{2\sinh(y/2)}
  +\frac{2\coth(y/2)}{\varepsilon-2x-y}\\[-1pt]
 &{}+\frac54\coth(x/2)+\frac1x .
\end{aligned}
\label{eq:pg-allocation}
\end{equation}
\end{proposition}
\begin{proof}
Substituting $\gamma=\tanh(y/2)$ and $\varepsilon_\ell=(\varepsilon-2x-y)/2$ into $B_A+\tau D$, removing constants, and dividing by $D-1$ gives Eq.~\eqref{eq:pg-allocation}.
The first and last terms are strictly convex in $y$ and $x$, respectively. The third term is convex.
The second term is log-convex because $[\log\coth(y/2)]''=\cosh(y)/\sinh^2(y)>0$ and $-\log(\varepsilon-2x-y)$ is convex.
Strict convexity and divergence at the boundary of the  domain give a unique interior minimizer. Its upper bound is at most that of the feasible split $\varepsilon_d=\varepsilon_r=\varepsilon_\ell=\varepsilon/5$, hence is $O(n)$ for fixed $\varepsilon,\tau$.
\end{proof}
Before Round~1, numerically minimize the two-variable objective in Eq.~\eqref{eq:pg-allocation} using only the public $n$ and $\varepsilon$.
Use $\varepsilon_d=x^\star$ per degree report, $\varepsilon_r=y^\star$ per assigned RR report, and $\varepsilon_\ell=(\varepsilon-2x^\star-y^\star)/2$ per relational report.
% The optimal bound is no larger than for epsilon/5 per mechanism and therefore attains the established minimax order.
% Fix all budgets before Round 1. Both endpoints consume degree and relational budgets; each edge has one RR report from its assigned endpoint.
